% \begin{tabularx}{\textwidth}[t]{|X|X|}
%   \hline \textbf{Notions et contenus} & \textbf{Capacités exigibles} \\
%   \hline \multicolumn{2}{|l|}{ 4.1. Changements de référentiel } \\
%   \hline 
%     Référentiel en translation rectiligne uniforme 
%     par rapport à un autre : transformation de 
%     Galilée, composition des vitesses.
%    & 
%     Relier la transformation de Galilée et la 
%     formule de composition des vitesses à la 
%     relation de Chasles et au caractère 
%     supposé absolu du temps.
%   \\
%   \hline 
%     Composition des vitesses et des 
%     accélérations dans le cas d'un référentiel en 
%     translation par rapport à un autre : point 
%     coïncident, vitesse d'entraînement, 
%     accélération d'entraînement.
%    & 
%     Exprimer la vitesse d'entraînement et 
%     l'accélération d'entraînement.
%   \\
%   \hline 
%     Composition des vitesses et des accélérations 
%     dans le cas d'un référentiel en rotation 
%     uniforme autour d'un axe fixe : point 
%     coïncident, vitesse d'entraînement, 
%     accélération d'entraînement, accélération de 
%     Coriolis.
%    & 
%     Exprimer la vitesse d'entraînement et 
%     l'accélération d'entraînement. 
%     Citer et utiliser l'expression de l'accélération 
%     de Coriolis.
%   \\
%   \hline
% \end{tabularx}

\vspace{0.25cm}
Dans la partie « Dynamique dans un référentiel non galiléen », l'étude du champ de pesanteur est conduite en supposant le référentiel géocentrique galiléen. De nombreuses applications permettent d'illustrer cette partie : le pendule de Foucault, la déviation vers l'est, les vents géostrophiques, les courants marins ; l'étude statique des marées constitue également une ouverture pertinente.

\begin{tabularx}{\textwidth}[t]{|X|X|}
  \hline 
  \textbf{Notions et contenus} & \textbf{Capacités exigibles} \\
  \hline 
  \multicolumn{2}{|l|}{4.1. Changements de référentiel} \\
  \hline 
  Champ de pesanteur terrestre : définition, évolution qualitative avec la latitude, ordres de grandeur. & Distinguer le champ de pesanteur et le champ gravitationnel. \\
  \hline 
  Capacité numérique : à l'aide d'un langage de programmation, illustrer un effet lié au caractère non galiléen du référentiel terrestre & \\
  \hline 
\end{tabularx}
